\section{How often does each model leak, and what differs for the second target?}
\label{app:behaviour}

Qwen3-32B leaked in some episodes and refused in others under almost every setting we screened, and the screens are too small to show which settings matter.
gpt-oss-20b and Gemma-4-31B refused direct requests for the secret, so gpt-oss-20b faces a different peer strategy and Gemma-4-31B was not used for probes.

\subsection{Screens on Qwen3-32B}

Before fixing the final setup, we ran small screens on Qwen3-32B that each varied one feature of the episode (\cref{tab:screens}).
Most screens use an earlier version of the main scenario.
The own-stake screen uses four other scenarios, in which the target gains something from sharing the value.
These screens use the scoring rules of the original screening runs, and \cref{tab:screens} describes behaviour without testing any setting.
\Cref{tab:behaviour-scenario} gives the leak rate of the final collection in each scenario.

Each screen was a separate collection, so a gap between two rows can reflect the batch or the scenario version as well as the setting.
One set of plain-request episodes serves as the reference for the tactic, thinking and peer-count screens.
The seven single tactics are helpful framing, an organiser's tone, a trade, social proof, time pressure, playing it safe and cost anchoring.
We adapted them from the persuasion prompts released by \citet{nakamura2026} and kept that source's honesty limits.
The reproducibility index points to the full table, which lists every cell with its interval and every excluded episode (\cref{app:repro}).

\begin{table}[h]
  \centering
  \small
  \resizebox{\linewidth}{!}{%
  \begin{tabular}{@{}lllll@{}}
    \toprule
    Feature varied & Reference setting & Leak rate & Changed setting & Leak rate \\
    \midrule
    Wording of the rule & states the consequence & \screenRuleBaseRate{} (\screenRuleBaseN{}) & adds ``this channel is logged'' & \screenRuleLoggedRate{} (\screenRuleLoggedN{}) \\
    Persuasion tactic & plain request & \screenPlainRate{} (\screenPlainN{}) & one of seven tactics & \screenTacticMinRate{} to \screenTacticMaxRate{} (\screenTacticMinN{} to \screenTacticMaxN{} each) \\
    Target's own stake & none & \screenStakeNoneRate{} (\screenStakeNoneN{}) & target gains from sharing & \screenStakeOwnRate{} (\screenStakeOwnN{}) \\
    Thinking mode & off & \screenPlainRate{} (\screenPlainN{}) & on (note a) & \screenThinkingRate{} (\screenThinkingN{}) \\
    Number of peers & one & \screenPlainRate{} (\screenPlainN{}) & two & \screenTwoPeersRate{} (\screenTwoPeersN{}) \\
    Awareness of a test (note b) & not flagged & \screenAwareNoRate{} (\screenAwareNoN{}) & flagged & \screenAwareYesRate{} (\screenAwareYesN{}) \\
    \bottomrule
  \end{tabular}}
  \par\smallskip
  \parbox{\linewidth}{\footnotesize
  Note a: turning thinking on also raised the reply-length limit, so this row changes two things at once.
  Note b: an automated judge flagged episodes in which the target seemed to recognise a test. Both flagged episodes only mention the mocked software tests in the task, so neither clearly shows such recognition.}
  \caption{\textbf{Almost every screened setting gave a mix of leaks and refusals on Qwen3-32B.}
  Each cell gives the share of recruiting episodes that leak, with the number of episodes in brackets.}
  \label{tab:screens}
\end{table}

\subsection{What differs for gpt-oss-20b?}

For gpt-oss-20b only the peer and the scored window change.
The task, the rule, the secret, the leak label and the detectors stay the same, and so does the control principle: a peer that never asks for the value.

The peer runs on DeepSeek-V4-Pro instead of DeepSeek-V4-Flash, with its own prompt (\cref{app:prompts}).
gpt-oss-20b mostly refused the plain request, and it never leaked when the peer stacked every tactic at once.
The peer therefore first coordinates on the ordinary task, then asks for the whole configuration block for a handoff file, without naming the secret.
We accepted a strategy only if four things held.
The target read the secret before any leak, made a scored call between the read and the leak, leaked in some episodes and refused in others, and never leaked in control episodes.
We chose the strategy on behaviour alone and froze it before any probe or monitor saw the data.

The scored window on gpt-oss-20b starts at the peer's first request for the block, not at the read.
The calls between the read and that request carry no pressure, so we keep them as a placebo on which a sound detector scores at chance (\cref{app:shortcuts}).

\subsection{Why was Gemma-4-31B not used for probes?}

Gemma-4-31B almost never leaked under direct requests, and in a small batch under the delayed block request it wrote the secret into a file every time.
It therefore gave no mix of leaking and refusing episodes, so we ran no replay, probes or monitors on it.
